%!TEX program = xelatex
% =====================================================================
%  PLANTILLA DE PRÁCTICAS - Escuela de Ingeniería Informática (ULPGC)
%  Compilar con XeLaTeX o LuaLaTeX (necesario para fontspec).
%  Imágenes esperadas en:  images/...
%    - images/logo_eii.png   -> logotipo ULPGC + EII de la portada
% =====================================================================
\documentclass[12pt,a4paper]{article}

% ---------------------------------------------------------------------
%  DATOS DEL DOCUMENTO  (lo único que hay que tocar en cada práctica)
% ---------------------------------------------------------------------
\newcommand{\autor}{Iván Pérez Díaz}
\newcommand{\fecha}{31/10/2026}
\newcommand{\asignatura}{Sistemas Robóticos Autónomos}
\newcommand{\numpractica}{Ejercicio Individual 1}
\newcommand{\titulopractica}{Respuesta Temporal.}

% Posición del logo de la portada: cuánto se "saca" del margen normal.
% Margen izquierdo real = 2.5 cm - \logodesplazax ; margen superior real = 2.8 cm - \logodesplazay
\newlength{\logodesplazax}\setlength{\logodesplazax}{1.2cm}
\newlength{\logodesplazay}\setlength{\logodesplazay}{1.8cm}
\newcommand{\logoancho}{0.73\textwidth}

% ---------------------------------------------------------------------
%  PAQUETES
% ---------------------------------------------------------------------
\usepackage{fontspec}
\usepackage[spanish,es-tabla]{babel}
\usepackage{geometry}
\usepackage{xcolor}
\usepackage{graphicx}
\usepackage{float}
\usepackage{caption}
\usepackage{setspace}
\usepackage{titlesec}
\usepackage{tocloft}
\usepackage{fancyhdr}
\usepackage{lastpage}
\usepackage[normalem]{ulem}
\usepackage{enumitem}
\usepackage{listings}
\usepackage{microtype}
\usepackage{array}
\usepackage{booktabs}
\usepackage[colorlinks=true,linkcolor=black,urlcolor=linkblue,citecolor=black]{hyperref}

% ---------------------------------------------------------------------
%  TIPOGRAFÍA  (Rubik si está instalada; si no, fuentes de respaldo)
% ---------------------------------------------------------------------
%  Texto normal : Rubik Regular
%  Títulos      : Rubik SemiBold  (\rubiksemibold)
\IfFontExistsTF{Rubik}{%
  \setmainfont{Rubik}[UprightFont=*-Regular,BoldFont=*-Bold,ItalicFont=*-Italic,BoldItalicFont=*-BoldItalic]%
  \setsansfont{Rubik}[UprightFont=*-Regular,BoldFont=*-Bold,ItalicFont=*-Italic,BoldItalicFont=*-BoldItalic]%
  \IfFontExistsTF{Rubik-SemiBold}{%
    \newfontfamily\rubiksemibold{Rubik-SemiBold}[BoldFont=Rubik-SemiBold]%
  }{%
    \newfontfamily\rubiksemibold{Rubik}[UprightFont=*-Bold]%
  }%
}{%
  \setmainfont{Latin Modern Sans}%
  \setsansfont{Latin Modern Sans}%
  \newfontfamily\rubiksemibold{Latin Modern Sans}[BoldFont=*-Bold]%
}
% Fuente monoespaciada del código: Consolas -> Fira Code -> DejaVu Sans Mono -> Latin Modern Mono
\IfFontExistsTF{Consolas}{\setmonofont{Consolas}}{%
  \IfFontExistsTF{Fira Code}{\setmonofont{Fira Code}[Scale=0.85]}{%
    \IfFontExistsTF{DejaVu Sans Mono}{\setmonofont{DejaVu Sans Mono}[Scale=0.82]}{%
      \setmonofont{Latin Modern Mono}}}}

% ---------------------------------------------------------------------
%  COLORES
% ---------------------------------------------------------------------
\definecolor{titleblue}{HTML}{0F4C97}     % título de portada
\definecolor{sectionblue}{HTML}{0B6CB5}   % secciones y subsecciones
\definecolor{linkblue}{HTML}{1155CC}      % enlaces
\definecolor{rulegray}{HTML}{B5B5B5}      % línea bajo los títulos
\definecolor{codebg}{HTML}{F5F5F5}        % fondo de los bloques de código
\definecolor{bgcmd}{HTML}{F4F4F4}         % fondo del código en línea

% ---------------------------------------------------------------------
%  PÁGINA, PÁRRAFOS E INTERLINEADO
% ---------------------------------------------------------------------
\geometry{a4paper,left=2.5cm,right=2.5cm,top=2.8cm,bottom=2.6cm,
          headheight=15pt,headsep=0.8cm,footskip=1.1cm}
\linespread{1.2}
\setlength{\parindent}{1.25cm}
\setlength{\parskip}{11pt}
\setlength{\emergencystretch}{2em}

% ---------------------------------------------------------------------
%  CABECERA Y PIE
% ---------------------------------------------------------------------
\pagestyle{fancy}
\fancyhf{}
% Zona superior: Rubik normal 12 pt, subrayada de un lado a otro (línea a todo el ancho)
\newcommand{\txtcabecera}{\fontsize{12}{14}\selectfont\sffamily}
\fancyhead[L]{\txtcabecera\autor}
\fancyhead[R]{\txtcabecera\asignatura\ - \numpractica}
\renewcommand{\headrulewidth}{0.6pt}
\fancyfoot[C]{\txtcabecera\thepage\ - \pageref*{LastPage}}
\fancypagestyle{plain}{\fancyhf{}%
  \fancyhead[L]{\txtcabecera\autor}%
  \fancyhead[R]{\txtcabecera\asignatura\ - \numpractica}%
  \fancyfoot[C]{\txtcabecera\thepage\ - \pageref*{LastPage}}%
  \renewcommand{\headrulewidth}{0.6pt}}

% ---------------------------------------------------------------------
%  TÍTULOS DE SECCIÓN (cada sección empieza en página nueva)
% ---------------------------------------------------------------------
\renewcommand{\thesection}{\arabic{section}}
\renewcommand{\thesubsection}{\arabic{section}.\arabic{subsection}}

% Título: Rubik semibold 20 pt
\titleformat{\section}
  {\rubiksemibold\fontsize{20}{26}\selectfont\color{sectionblue}}
  {\thesection.}{0.5em}{}
  [\vspace{0.4em}{\color{rulegray}\titlerule[0.6pt]}]
\titlespacing*{\section}{0pt}{0pt}{1.6em}

% Subtítulo: Rubik normal 16 pt, subrayado del mismo color que el texto
\newcommand{\subsecul}[1]{\uline{\thesubsection.\ #1}}
\titleformat{\subsection}
  {\sffamily\fontsize{16}{20}\selectfont\color{sectionblue}}
  {}{0pt}{\subsecul}
\titlespacing*{\subsection}{0pt}{1.6em}{1em}

\titleformat{\subsubsection}
  {\rubiksemibold\fontsize{14}{18}\selectfont\color{sectionblue}}
  {\thesubsubsection.}{0.5em}{}
\titlespacing*{\subsubsection}{0pt}{1.2em}{0.6em}

\let\oldsection\section
\renewcommand{\section}{\clearpage\oldsection}

% ---------------------------------------------------------------------
%  ÍNDICE
% ---------------------------------------------------------------------
\addto\captionsspanish{%
  \renewcommand{\contentsname}{\texorpdfstring{\sffamily\bfseries\fontsize{20}{26}\selectfont\color{sectionblue}\uline{Índice de Contenidos}}{Índice de Contenidos}}%
}
\renewcommand{\cfttoctitlefont}{\sffamily\bfseries\fontsize{20}{26}\selectfont\color{sectionblue}}
\renewcommand{\cftsecfont}{\sffamily\bfseries}
\renewcommand{\cftsecpagefont}{\sffamily\bfseries}
\renewcommand{\cftsubsecfont}{\sffamily}
\renewcommand{\cftsubsecpagefont}{\sffamily}
\renewcommand{\cftsecleader}{\cftdotfill{\cftdotsep}}
\renewcommand{\cftsubsecleader}{\cftdotfill{\cftdotsep}}
\setlength{\cftbeforesecskip}{0.6em}
\setlength{\cftbeforesubsecskip}{0.05em}
\setlength{\cftsecindent}{0pt}
\setlength{\cftsubsecindent}{0em}
\setlength{\cftsecnumwidth}{1.6em}
\setlength{\cftsubsecnumwidth}{2.6em}
\setlength{\cftaftertoctitleskip}{1em}
% ---------------------------------------------------------------------
%  FIGURAS Y PIES DE FIGURA  →  Figura 1. Texto del pie.
% ---------------------------------------------------------------------
\captionsetup{format=plain,labelsep=period,font={it,small},
              justification=centering,singlelinecheck=false,
              skip=6pt}
\captionsetup[figure]{justification=centering}
 
% \figura[ancho]{archivo}{etiqueta}{pie de figura}
%   Ej.:  \figura[0.8]{figura1.png}{fig:apache}{Configuración de Apache.}
\newcommand{\figura}[4][0.95]{%
  \begin{figure}[H]
    \centering
    \IfFileExists{images/#2}{%
      \includegraphics[width=#1\textwidth]{images/#2}}{%
      \fbox{\parbox{#1\textwidth}{\centering\sffamily\vspace{2em}
        Falta la imagen \texttt{images/#2}\vspace{2em}}}}
    \caption{#4}\label{#3}
  \end{figure}}
 
% Referencia a una figura: negrita + cursiva + subrayado + enlace
\newcommand{\figref}[1]{\hyperref[#1]{\textbf{\textit{\uline{Figura~\ref*{#1}}}}}}

% ---------------------------------------------------------------------
%  CÓDIGO
% ---------------------------------------------------------------------
% Código en línea: texto monoespaciado sobre fondo gris claro
\newcommand{\codigo}[1]{\colorbox{bgcmd}{\texttt{#1}}}

% Bloques de código: fondo gris claro, recuadro gris, numeración a la izquierda
\lstdefinestyle{codigo}{
  basicstyle=\ttfamily\fontsize{10}{12}\selectfont,
  backgroundcolor=\color{codebg},
  frame=single,             % Recuadro
  aboveskip=1.5em,          % Espacio extra ANTES del código
  rulecolor=\color{gray},   % Color del borde
  xleftmargin=16pt,         % Empuja todo el bloque hacia la derecha
  framexrightmargin=5pt,    % Aire en el interior derecho
  xrightmargin=5pt,         % Equilibra el margen derecho
  breaklines=true,
  captionpos=b,             % Títulos de código abajo
  numbers=left,
  numberstyle=\tiny\color{gray},
  keywordstyle=\color{blue},
  commentstyle=\color{gray},
  stringstyle=\color{teal},
  tabsize=2,
  showstringspaces=false
}
\lstset{style=codigo}

% ---------------------------------------------------------------------
% Definición de lenguajes extra
% ---------------------------------------------------------------------
\lstdefinelanguage{PowerShell}{
  morekeywords={ss, telnet, dig, ping, ssh},
  sensitive=false,
  morecomment=[l]{\#},
  morestring=[b]",
}
\lstdefinelanguage{Python}{
  keywords={def, return, if, else, elif, for, while, import, from, as, class, try, except},
  sensitive=true,
  morecomment=[l]{\#},
  morestring=[b]",
}

% Listas
\setlist{itemsep=0.2em,topsep=0.3em}

% =====================================================================
%  DOCUMENTO
% =====================================================================
\begin{document}

% ---------------------------------------------------------------------
%  PORTADA
% ---------------------------------------------------------------------
\begin{titlepage}
  \thispagestyle{empty}
  \vspace*{-\logodesplazay}%
  \noindent\hspace*{-\logodesplazax}%
  \IfFileExists{images/logo_eii.png}{%
    \includegraphics[width=\logoancho]{images/logo_eii.png}}{%
    \fbox{\parbox{0.5\textwidth}{\centering\sffamily Falta \texttt{images/logo\_eii.png}}}}

  \vspace*{\fill}
  \begin{center}
    \vspace{6em}
    \sffamily\bfseries\color{titleblue}
    {\fontsize{26}{32}\selectfont \numpractica: \titulopractica\par}
    \vspace{0.6em}
    \sffamily\bfseries\color{titleblue}
    {\Large (\asignatura)\par}
  \end{center}
  \vspace*{\fill}
  \vspace*{\fill}

  \noindent\sffamily\large \autor\hfill \fecha
\end{titlepage}

\pagenumbering{arabic}   % la portada no cuenta: el índice es la página 1

% ---------------------------------------------------------------------
%  ÍNDICE
% ---------------------------------------------------------------------
\tableofcontents
\clearpage

% ---------------------------------------------------------------------
%  CONTENIDO
% ---------------------------------------------------------------------
\section{Fichero .m para la Solución del Ejercicio}

\subsection{Valores de Entrada del Fichero}

El ejercicio propuesto consiste en la creación de un fichero .m que reciba como entrada los valores de dicha señal y los
correspondientes instantes de tiempo, y devuelva los siguientes resultados:

\begin{itemize}[label=$\bullet$]
  \item \textbf{Frecuencia amortiguada [Hz]:} Calculada como la inversa del pseudo-periodo de la respuesta.
  \item \textbf{Tiempo de establecimiento [seg]:} Calculado como el tiempo que tarda la señal en entrar en +-5\% de su valor final.
  \item \textbf{Tiempo de pico [seg]:} Instante de tiempo en el que se produce la primera sobreoscilación.
  \item \textbf{Error en régimen permanente}
\end{itemize}

Para la obtención de los resultados, se ha realizado el siguiente código en MatLab, fichero que se puede ver adjunto a la entrega o a través del siguiente enlace:
\href{https://github.com/ivanperezdiaz829/SRA/blob/main/Ejercicio_Individual_1/sol_ejercicio_1.m}{solución-ejercicio-1-Iván-Pérez-Díaz.m}.

El objetivo final es simular un sistema similar al de la \figref{fig:final-state}.

\figura[0.9]{final-state.png}{fig:final-state}{Estado final del sistema deseado.}

\subsection{Cabecera y Entradas}

\begin{lstlisting}[language=Matlab]
function [frec_amortiguada, temp_establecimiento, temp_pico, error_permanente] = sol_ejercicio_1(t, y, r)
\end{lstlisting}

La función recibe la señal \textbf{y} y sus correspondientes instantes de tiempo \textbf{t} y devuelve los cuatro resultados
pedidos en el enunciado del ejercicio, además, como señal de referencia, se ha añadido la entrada \textbf{r} que funciona como referencia del escalón y sirve para calcular el error en régimen permanente.

\begin{lstlisting}[language=Matlab]
if nargin < 3, r = 1; end
t = t(:); y = y(:);  
\end{lstlisting}

\begin{itemize}[label=$\bullet$]
  \item \codigo{nargin}\textbf{:} Es el número de argumentos con los que se ha llamado a la función. Si no se pasa \codigo{r} se asume un escalón unitario, encajando con la figura en donde la referencia es 1.
  \item \codigo{t(:)} \textbf{e} \codigo{y(:)}\textbf{:} Fuerzan a que ambos sean vectores columna, independientemente de cómo se pasen a la función. Arreglando posibles problemas cuando se use el bloque \codigo{To Workspace} de Simulink, que puede devolver vectores fila. 
\end{itemize}

\subsection{Valor Final}

\begin{lstlisting}[language=Matlab]
y_final = y(end);
\end{lstlisting}

Se toma como valor final la última muestra de la señal ya que en una respuesta subamortiguada las oscilaciones se atenuán y la señal se estabiliza, loq ue supone que la simulación dura lo suficiente para que ya esté asentada.
Es el \textbf{\textit{"valor final"}} que aparece en la \figref{fig:final-state}.

\subsection{Tiempo de Pico}

Para la detección de máximos locales, se ha programado la siguiente línea de código:

\begin{lstlisting}[language=Matlab]
idx_picos = find(diff(sign(diff(y))) < 0) + 1;
\end{lstlisting}

\begin{itemize}[label=$\bullet$]
  \item \codigo{diff(y)}\textbf{:} Calcula la diferencia entre muestras consecutivas, es decir, la dervidada discreta. Es positiva si la señal sube y negativa si baja.
  \item \codigo{sign(...)}\textbf{:} Convierte la dereviada discreta en 1 (sube), -1 (baja) o 0 (plana).
  \item \codigo{diff(sign(...))}\textbf{:} Calcula la diferencia entre muestras consecutivas de la derivada discreta. Si la señal pasa de subir a bajar, la diferencia es negativa, y si pasa de bajar a subir, es positiva.
  \item \codigo{find(... < 0)}\textbf{:} Devuelve las posiciones donde ocurre un cambio de subida a bajada, es decir, los máximos locales.
  \item El \codigo{+ 1} es necesario porque la función \codigo{diff} devuelve un vector con una longitud menor que el original, por lo que hay que compensar el índice. La posición \codigo{k} del resultado corresponde con la muestra \codigo{k+1} de \codigo{y}.
\end{itemize}

El resultado \codigo{idx\_picos}, contiene los índices de todos los picos, en orden temporal. El primero es la sobreoscilación principal y el segundo es el siguiente máximo local, mucho más pequeño.

En cuanto a la obtención del tiempo de pico, se ha programado la siguiente línea de código:

\begin{lstlisting}[language=Matlab]
if ~isempty(idx_picos)
  temp_pico = t(idx_picos(1));
else
  [~, idx] = max(y); temp_pico = t(idx);
end
\end{lstlisting}

La rama \codigo{else} es una protección. Si no se detecta ningún máximo local, se toma el máximo global con \codigo{max(y)}. La virgulilla descarta el valor máximo porque solo interesa su índice.

\subsection{Frecuencia Amortiguada}

\begin{lstlisting}[language=Matlab]
if numel(idx_picos) >= 2
  T_d = t(idx_picos(2)) - t(idx_picos(1));
  frec_amortiguada = 1 / T_d;
else
  frec_amortiguada = 0;
end
\end{lstlisting}

El enunciado define la frecuencia amortiguada como la inversa del pseudo-periodo. El pseudo-periodo \codigo{T\_d} es el tiempo entre dos picos consecutivos, y el código usa los dos primeros.
Su inversa está en Hz.

Se exige \codigo{numel(idx\_picos) >= 2} porque si hay menos de dos picos, no se puede calcular el pseudo-periodo. En ese caso, se devuelve 0.

\subsection{Tiempo de Establecimiento}

\begin{lstlisting}[language=Matlab]
banda = 0.05 * abs(y_final);
fuera_banda = find(abs(y - y_final) > banda);
if isempty(fuera_banda)
  temp_establecimiento = t(1);
else
  temp_establecimiento = t(min(fuera_banda(end) + 1, numel(t)));
end 
\end{lstlisting}

El enunciado define el tiempo de establecimiento como el tiempo que tarda la señal en entrar en +-5\% de su valor final y quedarse dentro. En la \figref{fig:final-state} son las líneas discontinuas alrededor del valor final.

\begin{itemize}[label=$\bullet$]
  \item \codigo{banda}\textbf{:} Es la semianchura de la franja de +-5\% del valor final. El \codigo{abs} evita una banda negativa en el caso de que el valor final fuese negativo.
  \item \codigo{abs(y - y\_final) > banda}\textbf{:} Marca las muestras que están fuera de la franja, y el \codigo{find} devuelve sus índices.
  \item \codigo{fuera\_banda(end)}\textbf{:} Es la última muestra que está fuera de al franja. Se usa la última y no la primera porque con oscilaciones, la señal puede entrar y salir de la banda varias veces anrtes de asentarse.
  \item \codigo{min(..., numel(t))}\textbf{:} Evita salirse del vector si la última muestra aún está fuera de la banda. En ese caso, devuelve el último instante.
\end{itemize}

\subsection{Error en Régimen Permanente}

\begin{lstlisting}[language=Matlab]
error_permanente = abs(r - y_final);
\end{lstlisting}

\section{Modelo y Simulación en Simulink}

\subsection{Diagrama Completo}

Para crear el diagrama, tras guardar el proyecto se ha de ir a \textbf{Simulink} y poner los siguientes bloques para poder ejecutar correctamente la simulación y poder ver la onda final:

\begin{itemize}[label=$\bullet$]
  \item \textbf{Step:} Simula la entrada escalón (\codigo{r}).
  \item \textbf{Transfer Fcn:} Es la simulación del sistema subamortiguado.
  \item \textbf{To Workspace:} Es el espacio de trabajo hecho en MatLab explicado en el \textbf{primer punto} de este documento.
  \item \textbf{Scope:} Se usa para ver la señal resultante.
\end{itemize}

Una vez puestos todos los bloques, el diagrama se ve tal y como se aprecia en la \figref{fig:diagrama}.

\figura[0.9]{diagrama.png}{fig:diagrama}{Diagrama en Simulink.}

\subsection{Bloque Step}

\figura[0.49]{step.png}{fig:step}{Parámetros del Bloque Step.}

\subsection{Bloque Transfer Fcn}

\figura[0.48]{transfer-fcn.png}{fig:fcn}{Parámetros del Bloque Transfer Fcn.}

\subsection{Bloque To workspace}

\figura[0.6]{toWorkspace.png}{fig:workspace}{Parámetros del bloque toWorkspace.}

\section{Script de prueba del fichero .m.}

\subsection{Objetivo del script.}

El script comprueba que la función \codigo{sol\_ejercicio\_1} funciona bien. Para ello crea una señal de la que se conocen los resultados exactos, se la pasa a la función y compara lo que devuelve con los valores teóricos. Por último, dibuja la señal marcando cada resultado.

La señal es la misma que produce el modelo de Simulink de la sección anterior ($\omega_n = 20$, $\zeta = 0.3$ y $K = 0.9$), pero aquí se genera con su fórmula. Para ver el código, véase el documento: \href{https://github.com/ivanperezdiaz829/SRA/blob/main/Ejercicio_Individual_1/Ejercicio_1_test.m}{Ejercicio-1-Test.m}

\subsection{Parámetros del Sistema}

\begin{lstlisting}[language=Matlab]
r    = 1;        % Referencia del escalon
K    = 0.9;      % Ganancia estatica (valor final = K*r)
wn   = 20;       % Frecuencia natural [rad/s]
zeta = 0.3;      % Amortiguamiento (0 < zeta < 1)
Tm    = 1e-3;    % Paso de muestreo [s]
t_fin = 2;       % Duracion [s]
\end{lstlisting}

Se definen los datos del sistema. Como \codigo{zeta} está entre 0 y 1, el sistema es subamortiguado. La ganancia \codigo{K} se ha puesto distinta de 1 a propósito: así la señal se estabiliza en 0.9 y no en la referencia 1, lo que permite comprobar que el error permanente se calcula bien. El paso de 1 ms y los 2 s de duración son suficientes para que la señal acabe asentada.

\subsection{Generación de la Señal}

\begin{lstlisting}[language=Matlab]
t   = (0:Tm:t_fin)';
wd  = wn*sqrt(1 - zeta^2);
phi = acos(zeta);
y   = K*r*(1 - exp(-zeta*wn*t)/sqrt(1 - zeta^2) .* sin(wd*t + phi));
\end{lstlisting}

Se crea el vector de tiempos y se calcula la señal con la fórmula de la respuesta de un sistema de segundo orden subamortiguado ante un escalón. La exponencial hace que las oscilaciones se vayan atenuando, y el seno es la oscilación, que ocurre a la frecuencia amortiguada \codigo{wd}.

\subsection{Llamada a la Función y Valores Teóricos}

\begin{lstlisting}[language=Matlab]
[frec_amortiguada, temp_establecimiento, temp_pico, error_permanente] = ...
    sol_ejercicio_1(t, y, r);

fd_teo  = wd/(2*pi);
tp_teo  = pi/wd;
ts_teo  = -log(0.05*sqrt(1 - zeta^2))/(zeta*wn);
ess_teo = abs(r - K*r);
\end{lstlisting}

Primero se llama a la función con la señal, el tiempo y la referencia. Después se calculan los valores que \textbf{deberían} salir, usando las fórmulas del sistema de segundo orden:

\begin{itemize}[label=$\bullet$]
  \item \textbf{Frecuencia amortiguada:} $f_d = \omega_d / 2\pi$.
  \item \textbf{Tiempo de pico:} $t_p = \pi / \omega_d$, que es medio pseudo-periodo.
  \item \textbf{Tiempo de establecimiento:} Instante en el que las oscilaciones quedan por debajo del 5\%, $t_s = -\ln\left(0.05\sqrt{1-\zeta^2}\right) / (\zeta\omega_n)$.
  \item \textbf{Error permanente:} $e_{ss} = |r - Kr|$.
\end{itemize}

\subsection{Resultados por Pantalla}

\begin{lstlisting}[language=Matlab]
fprintf('\n=== RESULTADOS ===\n');
fprintf('%-30s %12s %12s\n', 'Magnitud', 'Calculado', 'Teorico');
fprintf('%-30s %12.4f %12.4f\n', 'Frecuencia amortiguada [Hz]', ...
        frec_amortiguada, fd_teo);
% (igual para el resto de magnitudes)
\end{lstlisting}

Se imprime una tabla con el valor calculado por la función y el teórico de cada magnitud. El resultado se puede ver en la \figref{fig:terminal}.

\figura[0.7]{terminal.png}{fig:terminal}{Resultados numéricos en la terminal.}

\subsection{Gráfica de la Señal}

\begin{lstlisting}[language=Matlab]
plot(t, y, 'b', 'LineWidth', 1.8);
plot([t(1) t(end)], [r r], 'k-');
plot([t(1) t(end)], [y_final y_final], 'k--');
plot([t(1) t(end)], 1.05*y_final*[1 1], 'b:');
plot([t(1) t(end)], 0.95*y_final*[1 1], 'b:');
plot(temp_pico, y_pico, 'ro', 'MarkerFaceColor', 'r');
plot([temp_establecimiento temp_establecimiento], [0 1.05*y_final], 'g--');
plot([t(end) t(end)], [y_final r], 'm-', 'LineWidth', 3);
\end{lstlisting}

Se dibuja la señal y se marca sobre ella cada resultado obtenido, como se ve en la \figref{fig:salida}:

\begin{itemize}[label=$\bullet$]
  \item \textbf{Señal} (azul): la respuesta $y_2(t)$.
  \item \textbf{Referencia y valor final:} la línea negra continua es la referencia (1) y la discontinua el valor donde se estabiliza la señal (0.9).
  \item \textbf{Banda del $\pm$5\%:} las dos líneas azules punteadas alrededor del valor final.
  \item \textbf{Tiempo de pico:} el punto rojo con su línea vertical.
  \item \textbf{Tiempo de establecimiento:} la línea vertical verde.
  \item \textbf{Error permanente:} el segmento magenta de la derecha, entre el valor final y la referencia.
\end{itemize}

Al final, el script guarda la figura en PDF y PNG para la entrega.

\figura[0.8]{diagrama-salida.png}{fig:salida}{Respuesta temporal del sistema.}

\subsection{Interpretación de los Resultados}

La señal de la \figref{fig:salida} es la típica de un sistema subamortiguado: sube rápido, se pasa del valor final (llega a un pico de aproximadamente 1.24), baja, y oscila cada vez menos hasta quedarse quieta. Fíjese en que se estabiliza en 0.9 y no en la referencia 1. El significado de cada resultado de la \figref{fig:terminal} es:

\begin{itemize}[label=$\bullet$]
  \item \textbf{Frecuencia amortiguada} (3.0395 Hz)\textbf{:} Son las oscilaciones por segundo de la señal, es decir, la inversa del tiempo entre dos picos (unos 0.33 s). Es algo menor que la frecuencia natural ($\omega_n/2\pi \approx 3.18$ Hz) porque el amortiguamiento frena las oscilaciones.
  \item \textbf{Tiempo de pico} (0.1650 s)\textbf{:} Es lo que tarda la señal en llegar a su primer máximo, el punto rojo. Es más o menos medio pseudo-periodo, porque la señal parte de 0 y tarda media oscilación en llegar arriba.
  \item \textbf{Tiempo de establecimiento} (0.5070 s)\textbf{:} Es el momento a partir del cual la señal ya no sale de la banda de $\pm$5\% alrededor del valor final (línea verde). Indica lo rápido que se estabiliza el sistema.
  \item \textbf{Error en régimen permanente} (0.1000)\textbf{:} Es la distancia entre la referencia (1) y el valor donde se estabiliza la señal (0.9), el segmento magenta. Significa que el sistema se estabiliza pero no llega exactamente a lo que se le pide, porque su ganancia es 0.9 y no 1.
\end{itemize}

\subsection{Comparación con los Valores Teóricos}

\begin{table}[H]
  \centering
  \begin{tabular}{lccc}
    \toprule
    \textbf{Magnitud} & \textbf{Calculado} & \textbf{Teórico} & \textbf{Error relativo} \\
    \midrule
    Frecuencia amortiguada [Hz]   & 3.0395 & 3.0365 & 0.10\% \\
    Tiempo de establecimiento [s] & 0.5070 & 0.5071 & 0.02\% \\
    Tiempo de pico [s]            & 0.1650 & 0.1647 & 0.18\% \\
    Error en régimen permanente   & 0.1000 & 0.1000 & 0\% \\
    \bottomrule
  \end{tabular}
  \caption{Comparación entre los valores calculados y los teóricos.}
\end{table}

Los resultados calculados coinciden con los teóricos con diferencias inferiores al 0.2\%. Esa pequeña diferencia se debe a que la función trabaja con muestras separadas 1 ms, así que solo puede detectar los instantes con esa precisión, mientras que los teóricos se calculan de forma exacta. Por tanto, la función \codigo{sol\_ejercicio\_1} calcula correctamente las cuatro magnitudes pedidas.

\section{Bibliografía}

\subsection{Uso de Inteligencia Artificial}

El código del script de prueba (\codigo{test\_sol\_ejercicio\_1.m}) ha sido programado íntegramente con \textbf{Claude.ai}. Esta misma herramienta se ha utilizado como apoyo en la revisión del formato del documento (realizado en LaTeX). El contenido ha sido revisado y los resultados verificados ejecutando el código en MatLab.

\begin{itemize}[label=$\bullet$]
  \item Anthropic. \textit{Claude.ai}. \href{https://claude.ai}{https://claude.ai}. Consultado el 31/10/2026.
\end{itemize}

\subsection{Documentación Oficial de MatLab}

Referencias a la documentación oficial de MathWorks para las funciones empleadas en el código. Consultadas el 31/10/2026.

\begin{itemize}[label=$\bullet$]
  \item MathWorks. \textit{diff}: Differences and approximate derivatives. \href{https://www.mathworks.com/help/matlab/ref/diff.html}{diff.html}
  \item MathWorks. \textit{sign}: Sign function. \href{https://www.mathworks.com/help/matlab/ref/sign.html}{sign.html}
  \item MathWorks. \textit{find}: Find indices and values of nonzero elements. \href{https://www.mathworks.com/help/matlab/ref/find.html}{find.html}
  \item MathWorks. \textit{max}: Maximum elements of an array. \href{https://www.mathworks.com/help/matlab/ref/max.html}{max.html}
  \item MathWorks. \textit{nargin}: Number of function input arguments. \href{https://www.mathworks.com/help/matlab/ref/nargin.html}{nargin.html}
  \item MathWorks. \textit{fprintf}: Write data to text file. \href{https://www.mathworks.com/help/matlab/ref/fprintf.html}{fprintf.html}
  \item MathWorks. \textit{plot}: 2-D line plot. \href{https://www.mathworks.com/help/matlab/ref/plot.html}{plot.html}
  \item MathWorks. \textit{yline}: Horizontal line with constant y-value. \href{https://www.mathworks.com/help/matlab/ref/yline.html}{yline.html}
  \item MathWorks. \textit{xline}: Vertical line with constant x-value. \href{https://www.mathworks.com/help/matlab/ref/xline.html}{xline.html}
  \item MathWorks. \textit{exportgraphics}: Save plot or graphics content to file. \href{https://www.mathworks.com/help/matlab/ref/exportgraphics.html}{exportgraphics.html}
\end{itemize}

\subsection{Documentación Oficial de Simulink}

Referencias a la documentación oficial de MathWorks para los bloques utilizados en el modelo. Consultadas el 31/10/2026.

\begin{itemize}[label=$\bullet$]
  \item MathWorks. \textit{Step}: Generate step function. \href{https://www.mathworks.com/help/simulink/slref/step.html}{step.html}
  \item MathWorks. \textit{Transfer Fcn}: Model linear system by transfer function. \href{https://www.mathworks.com/help/simulink/slref/transferfcn.html}{transferfcn.html}
  \item MathWorks. \textit{To Workspace}: Write data to workspace. \href{https://www.mathworks.com/help/simulink/slref/toworkspace.html}{toworkspace.html}
  \item MathWorks. \textit{Scope}: Display signals generated during simulation. \href{https://www.mathworks.com/help/simulink/slref/scope.html}{scope.html}
\end{itemize}

\subsection{Documentación de la Asignatura}

\begin{itemize}[label=$\bullet$]
  \item Universidad de Las Palmas de Gran Canaria. \textit{Sistemas Robóticos Autónomos}: enunciado del Ejercicio Individual 1 y material de la asignatura. Campus Virtual de la ULPGC (Moodle). \href{https://aep27.ulpgc.es/course/view.php?id=2176}{link}
\end{itemize}

\end{document}